\SourceLeaf{348}%
tive bei concav-convexen Statt hat. Wenn das in Rede stehende Stück der krummen Fläche in dieser Beziehung aus Theilen ungleicher Art besteht, so werden noch nähere Bestimmungen nothwendig, die hier übergangen werden müssen.

Die Vergleichung des Inhalts zweier einander correspondirender Stücke der krummen Fläche und der Oberfläche der Hülfskugel führt nun (auf dieselbe Art wie z.\,B. aus der Vergleichung von Volumen und Masse der Begriff von Dichtigkeit hervorgeht) zu einem neuen Begriffe. Der Verf. nennt nämlich \textit{Krümmungsmaass} in einem Punkt der krummen Fläche den Werth des Bruches, dessen Nenner der Inhalt eines unendlich kleinen Stücks der krummen Fläche in diesem Punkt, und der Zähler der Inhalt des entsprechenden Stücks der Fläche der Hülfskugel, oder die ganze Krümmung jenes Elements ist. Man sieht, dass, in dem Sinn des Verf., ganze Krümmung und Krümmungsmaass bei krummen Flächen dem analog ist, was bei krummen Linien resp. Amplitudo und schlechthin Krümmung genannt wird; er fand Bedenken, die letztern mehr durch Gewohnheit als wegen Angemessenheit recipirten Ausdrücke auf die krummen Flächen zu übertragen. Uebrigens liegt weniger an den Benennungen selbst, als daran, dass ihre Einführung durch prägnante Sätze gerechtfertigt wird.

Die Auflösung der Aufgabe, das Krümmungsmaass in jedem Punkt einer krummen Fläche zu finden, erscheint in verschiedener Gestalt, nach Maassgabe der Art, wie die Natur der krummen Fläche gegeben ist. Die einfachste Art ist, indem die Punkte im Raum allgemein durch drei rechtwinklige Coordinaten \(x,y.\ z\) unterschieden werden, eine Coordinate als Function der beiden andern darzustellen: dabei erhält man den einfachsten Ausdruck für das Krümmungsmaass. Zugleich ergibt sich aber ein merkwürdiger Zusammenhang zwischen diesem Krümmungsmaass und den Krümmungen derjenigen Curven, die durch den Schnitt der krummen Fläche mit Ebenen senkrecht auf dieselbe, hervorgehen. Bekanntlich hat \textsc{Euler} zuerst gezeigt, dass zwei dieser schneidenden Ebenen, die einander gleichfalls unter einem rechten Winkel schneiden, die Eigenschaft haben, dass in der einen der grösste, in der andern der kleinste Krümmungshalbmesser Statt findet, oder richtiger, dass in ihnen die beiden äussersten Krümmungen vorkommen. Hier ergibt sich nun aus dem erwähnten Ausdruck für das Krümmungsmaass, dass dieses einem Bruche gleich wird, dessen Zähler die Einheit, der Nenner das Product der beiden äussersten Krümmungshalbmesser wird. --- Weniger einfach wird der Ausdruck für das Krümmungsmaass, wenn
